% ===== PRÉAMBULE CANONIQUE — à coller VERBATIM en tête de CHAQUE .tex =====
\documentclass[11pt,a4paper]{article}
\usepackage[utf8]{inputenc}\usepackage[T1]{fontenc}\usepackage{lmodern}
\usepackage[french]{babel}
\usepackage{amsmath,amssymb,mathtools}\usepackage{icomma}
\usepackage{siunitx}\sisetup{locale=FR}  % \qty{}{}, \si{}, \ang{}
\usepackage{xcolor}\usepackage{colortbl}
\usepackage{tikz}\usepackage{pgfplots}\pgfplotsset{compat=1.18}
\usepackage[siunitx,europeanresistors]{circuitikz} % circuits (élec.)
\usepackage[most]{tcolorbox}\usepackage{enumitem}\usepackage{array,booktabs}
\usepackage[a4paper,margin=2.2cm]{geometry}
\usetikzlibrary{arrows.meta,calc,angles,quotes,positioning,patterns,%
  backgrounds,shapes.geometric,decorations.pathmorphing,decorations.markings,intersections}
\definecolor{primary}{RGB}{222,49,99}\definecolor{primarydark}{RGB}{150,22,55}
\definecolor{defcol}{RGB}{0,114,178}\definecolor{methcol}{RGB}{52,92,148}
\definecolor{excol}{RGB}{0,135,90}\definecolor{attncol}{RGB}{200,40,55}
\tcbset{boxbase/.style={enhanced,breakable,boxrule=0.9pt,arc=2.5pt,
  left=3mm,right=3mm,top=2.3mm,bottom=2.3mm,fonttitle=\bfseries,coltitle=white}}
% --- boîtes TUTORIEL (noms EXACTS, ne pas renommer) ---
\newtcolorbox{cle}[1][]{boxbase,colback=defcol!6,colframe=defcol,
  title={\textsc{À retenir}\ifx\relax#1\relax\relax\else\textmd{ : \itshape #1}\fi}}
\newtcolorbox{methode}[1][]{boxbase,colback=methcol!7,colframe=methcol,sharp corners,
  title={\textsc{Méthode}\ifx\relax#1\relax\relax\else\textmd{ --- \itshape #1}\fi}}
\newtcolorbox{exemplebox}[1][]{boxbase,colback=excol!5,colframe=excol,
  title={\textsc{Exemple}\ifx\relax#1\relax\relax\else\textmd{ : \itshape #1}\fi}}
\newtcolorbox{piege}[1][]{boxbase,colback=attncol!6,colframe=attncol,
  title={\textsc{Piège}\ifx\relax#1\relax\relax\else\textmd{ : \itshape #1}\fi}}
\newtcolorbox{remarque}[1][]{boxbase,colback=black!4,colframe=black!45,
  title={\textsc{Remarque}\ifx\relax#1\relax\relax\else\textmd{ : \itshape #1}\fi}}
% --- boîtes EXERCICE / CORRIGÉ (noms EXACTS, auto-comptées) ---
\newtcolorbox[auto counter]{exo}[1][]{boxbase,colback=primary!4,colframe=primary,
  title={\textsc{Exercice}~\thetcbcounter\ifx\relax#1\relax\relax\else\textmd{ --- \itshape #1}\fi}}
\newtcolorbox[auto counter]{corrige}[1][]{boxbase,colback=excol!4,colframe=excol,
  title={\textsc{Corrigé}~\thetcbcounter\ifx\relax#1\relax\relax\else\textmd{ --- \itshape #1}\fi}}
\newcommand{\vv}[1]{\vec{#1}}\newcommand{\ds}{\displaystyle}
\newcommand{\R}{\mathbb{R}}\newcommand{\N}{\mathbb{N}}\newcommand{\Z}{\mathbb{Z}}
% (un \usepackage{fancyhdr}+en-tête est possible ; il est ignoré côté web)
% =========================================================================
\begin{document}
\sisetup{per-mode=symbol}

\begin{corrige}[le choix de la référence est arbitraire]
\begin{enumerate}[label=\alph*)]
  \item Référence au sol : $E_p=mgh=4\times10\times3=\qty{120}{\joule}$.
  \item Référence sur la table : la hauteur au-dessus de la table est $3-1=\qty{2}{\metre}$, donc
        $E_p=4\times10\times2=\qty{80}{\joule}$.
  \item Après la descente, le drone est à \qty{1}{\metre} du sol.
        \begin{itemize}[leftmargin=1.4em,topsep=1pt]
          \item réf. sol : $\Delta E_p=mg(1)-mg(3)=40-120=\qty{-80}{\joule}$ ;
          \item réf. table : $\Delta E_p=mg(0)-mg(2)=0-80=\qty{-80}{\joule}$.
        \end{itemize}
        Les $E_p$ diffèrent selon la référence, mais la \emph{variation} $\Delta E_p=\qty{-80}{\joule}$
        est \textbf{la même} : seule $\Delta E_p$ a un sens physique.
\end{enumerate}
\end{corrige}

\begin{corrige}[l'énergie élastique varie comme le carré de l'élongation]
\begin{enumerate}[label=\alph*)]
  \item $x_1=\qty{0,04}{\metre}$ : $E_p=\tfrac12\times300\times(0,04)^2=\tfrac12\times300\times0,0016=\qty{0,24}{\joule}$.
  \item $x_2=\qty{0,08}{\metre}$ : $E_p=\tfrac12\times300\times(0,08)^2=\tfrac12\times300\times0,0064=\qty{0,96}{\joule}$.
        L'élongation a doublé, $E_p$ a été multipliée par $\dfrac{0,96}{0,24}=4$ (facteur $2^2$).
\end{enumerate}
\end{corrige}

\begin{corrige}[énergie stockée entre deux élongations]
Avec $x_1=\qty{0,02}{\metre}$ et $x_2=\qty{0,06}{\metre}$ :
\[
W=\tfrac12 k\,(x_2^2-x_1^2)=\tfrac12\times500\times\bigl((0,06)^2-(0,02)^2\bigr)
=250\times(0,0036-0,0004)=250\times0,0032=\qty{0,8}{\joule}.
\]
\end{corrige}

\begin{corrige}[influence de la masse et de la raideur sur la période]
\begin{enumerate}[label=\alph*)]
  \item $T\propto\sqrt{m}$ : si $m$ est multipliée par $4$, alors $T$ est multipliée par
        $\sqrt{4}=2$ (la période \emph{double}).
  \item $T\propto\dfrac{1}{\sqrt{k}}$ : si $k$ est multipliée par $4$, alors $T$ est multipliée par
        $\dfrac{1}{\sqrt{4}}=\dfrac12$ (la période est \emph{divisée par 2}).
\end{enumerate}
\end{corrige}

\begin{corrige}[l'énergie potentielle ne dépend que de la hauteur]
\begin{enumerate}[label=\alph*)]
  \item $\Delta E_p=mgh=20\times10\times3=\qty{600}{\joule}$.
  \item \textbf{Non}, $\Delta E_p=mgh$ ne dépend que de la hauteur $h$ (et de $m$, $g$), \emph{pas} du
        chemin suivi : la rampe (\qty{6}{\metre}) et la montée verticale (\qty{3}{\metre}) donnent la
        même variation d'énergie potentielle. Le poids est une force \emph{conservative}.
\end{enumerate}
\end{corrige}
\end{document}
